%% Chương 2: Cơ sở lý thuyết
\clearpage
\phantomsection
\section*{CHƯƠNG 2. CƠ SỞ LÝ THUYẾT}
\addcontentsline{toc}{section}{\numberline{}CHƯƠNG 2. CƠ SỞ LÝ THUYẾT}
\setcounter{section}{2}
\setcounter{subsection}{0}
\setcounter{figure}{0}
\setcounter{table}{0}
\setcounter{equation}{0}

\subsection{Chuẩn H.264/AVC}

\subsubsection{Lịch sử và vị trí của H.264}

H.264/AVC (Advanced Video Coding) là chuẩn nén video được phát triển chung bởi
ITU-T Video Coding Experts Group (VCEG) và ISO/IEC Moving Picture Experts Group (MPEG),
được công bố lần đầu năm 2003 và liên tục cập nhật đến nay \cite{itu_h264}.
Chuẩn này có tên gọi song song là H.264 (trong hệ thống ITU-T) và MPEG-4 AVC
(trong hệ thống ISO/IEC 14496-10).

So với các chuẩn tiền nhiệm như H.263 hay MPEG-2, H.264 đạt mức nén gấp 2 lần
tại cùng chất lượng hình ảnh. Điều này đạt được thông qua nhiều kỹ thuật tiên tiến
như dự đoán chuyển động đa tham chiếu (multi-reference motion prediction), biến đổi
integer DCT 4$\times$4, mã hoá entropy CABAC (Context-Adaptive Binary Arithmetic Coding),
và deblocking filter. Sự kết hợp giữa chất lượng nén cao và tính tương thích phần cứng
rộng làm cho H.264 trở thành chuẩn video phổ biến nhất trong suốt hơn hai thập kỷ.

\subsubsection{Kiến trúc phân lớp của H.264}

Chuẩn H.264 được thiết kế theo kiến trúc hai lớp tách biệt:

\begin{itemize}
  \item \textbf{VCL (Video Coding Layer):} Lớp mã hoá video thuần túy, chịu trách
        nhiệm nén dữ liệu ảnh. VCL xử lý các slice (đơn vị mã hoá độc lập),
        macroblock (khối ảnh 16$\times$16 pixel), và các phép biến đổi toán học.

  \item \textbf{NAL (Network Abstraction Layer):} Lớp trừu tượng mạng, đóng gói
        dữ liệu VCL thành các đơn vị NALU (NAL Unit) phù hợp với truyền tải
        qua các giao thức mạng khác nhau (RTP, HTTP, v.v.) hoặc lưu trữ.
\end{itemize}

Sự tách biệt này cho phép cùng một chuẩn nén hoạt động trên nhiều môi trường
truyền tải khác nhau mà không cần thay đổi phần mã hoá video cốt lõi.

\subsubsection{Cấu trúc NAL Unit (NALU)}

Mỗi NAL Unit là đơn vị cơ bản của luồng bit H.264. Trong định dạng Annex B
(byte stream format --- định dạng được sử dụng trong nghiên cứu này), mỗi NALU
được ngăn cách bởi start code:

\begin{itemize}
  \item Start code 4 bytes: \texttt{0x00 0x00 0x00 0x01}
  \item Hoặc start code 3 bytes: \texttt{0x00 0x00 0x01} (cho các NALU không phải đầu tiên trong access unit)
\end{itemize}

Sau start code là payload của NALU, bắt đầu bằng \textbf{NAL header} (1 byte) với cấu trúc:

\begin{table}[ht]
\centering
\caption{Cấu trúc byte NAL header (1 byte)}
\label{tab:nal_header}
\begin{tabular}{cccl}
\toprule
Bit & Tên trường & Kích thước & Ý nghĩa \\
\midrule
7   & forbidden\_zero\_bit  & 1 bit  & Luôn = 0 (kiểm tra lỗi) \\
5-6 & nal\_ref\_idc         & 2 bits & Mức độ tham chiếu (0: không tham chiếu) \\
0-4 & nal\_unit\_type       & 5 bits & Loại NALU (xem Bảng~\ref{tab:nal_types}) \\
\bottomrule
\end{tabular}
\end{table}

\subsubsection{Phân loại NAL Unit}

Trường \texttt{nal\_unit\_type} (5 bit thấp của NAL header) xác định loại dữ liệu
trong NALU:

\begin{table}[ht]
\centering
\caption{Các loại NAL Unit quan trọng trong H.264}
\label{tab:nal_types}
\begin{tabular}{cll}
\toprule
\texttt{nal\_unit\_type} & Tên & Ý nghĩa \\
\midrule
1 & Coded slice (non-IDR) & P-frame (predictive) hoặc B-frame (bi-predictive) \\
2 & Coded slice data partition A & Phân vùng dữ liệu (ít phổ biến) \\
5 & IDR (Instantaneous Decoding Refresh) slice & I-frame (Intra-coded) --- keyframe \\
6 & SEI                   & Supplemental Enhancement Information \\
7 & SPS                   & Sequence Parameter Set (thông số chung) \\
8 & PPS                   & Picture Parameter Set (thông số ảnh) \\
9 & Access Unit Delimiter & Dấu phân cách access unit \\
\bottomrule
\end{tabular}
\end{table}

\textbf{Phân loại trong nghiên cứu:}
\begin{itemize}
  \item \textbf{VCL NALU (Video Coding Layer):} Type 1 và Type 5 chứa dữ liệu ảnh
        thực sự và là đối tượng mã hoá trong nghiên cứu này.
  \item \textbf{Non-VCL NALU:} Type 7 (SPS), Type 8 (PPS), Type 6 (SEI) chứa
        thông số cấu hình và metadata --- \textbf{không bị mã hoá}.
\end{itemize}

Trong file thực nghiệm \texttt{output.h264}, phân phối NALU như sau:
Type 5 (IDR/I-frame): 155 NALU; Type 1 (P/B-frame): 18.882 NALU; các loại khác: 311 NALU.
Chiến lược S1 mã hoá 19.037 VCL NALU (Type 1 + Type 5, bỏ qua 311 non-VCL),
S2 chỉ mã hoá 155 NALU Type 5.

\subsubsection{Intra Prediction và hệ số DC}

Để hiểu tầm quan trọng của hệ số DC, cần nắm cơ chế dự đoán intra của H.264.
Với I-frame (IDR slice, Type 5), mọi macroblock được mã hoá bằng \textbf{intra prediction}
--- dự đoán khối ảnh từ các khối đã được mã hoá trong cùng frame, không tham chiếu
frame khác. Quá trình xử lý một macroblock luma (16$\times$16 pixel) diễn ra như sau:

\begin{enumerate}
  \item \textbf{Intra prediction:} Dự đoán giá trị pixel từ các pixel lân cận đã biết.
        H.264 hỗ trợ 13 chế độ dự đoán intra 4$\times$4 và 4 chế độ 16$\times$16.
        Kết quả là \emph{residual} (sai lệch giữa thực tế và dự đoán).

  \item \textbf{Biến đổi integer DCT:} Residual $R$ (4$\times$4 hoặc 8$\times$8) được
        biến đổi bằng integer approximation của DCT:
        \begin{equation}
        Y = H \cdot R \cdot H^T
        \end{equation}
        trong đó $H$ là ma trận biến đổi integer DCT. Kết quả $Y$ có 16 hệ số,
        trong đó \textbf{hệ số DC là $Y[0,0]$} --- hệ số tần số thấp nhất,
        biểu diễn giá trị trung bình của khối residual.

  \item \textbf{Lượng tử hoá:} Các hệ số được lượng tử hoá theo QP (quantization parameter).
        Hệ số DC thường có weight lượng tử hoá khác với AC để bảo toàn chất lượng.

  \item \textbf{Mã hoá entropy:} Các hệ số sau lượng tử hoá được mã hoá bằng CAVLC
        hoặc CABAC và ghép vào bitstream.
\end{enumerate}

\subsubsection{Vị trí hệ số DC trong luồng bit NALU}

Qua phân tích bitstream bằng công cụ \texttt{h264\_analyze} \cite{h264bitstream},
chúng tôi xác định rằng đối với file \texttt{output.h264} (H.264 Baseline Profile),
các hệ số DC của slice header chiếm vùng bytes từ offset 19 đến 43 trong payload
NALU VCL (sau start code):

\begin{itemize}
  \item \texttt{DC\_OFFSET = 19}: Offset tính từ đầu payload NALU (sau start code).
  \item \texttt{DC\_SIZE = 25}: Kích thước vùng DC, 25 bytes.
  \item Vị trí tuyệt đối: \texttt{nal\_pos + 19} đến \texttt{nal\_pos + 43}
        (với \texttt{nal\_pos} là vị trí byte đầu của payload NALU trong file).
\end{itemize}

\begin{figure}[ht]
\centering
%% TikZ: DC offset diagram — byte layout trong NALU VCL
\begin{tikzpicture}[
  box/.style={draw, minimum height=0.7cm, font=\small},
  header/.style={box, fill=gray!20, minimum width=1.8cm},
  dc/.style={box, fill=orange!40, minimum width=3.5cm},
  rest/.style={box, fill=blue!10, minimum width=2.5cm},
  startcode/.style={box, fill=red!15, minimum width=1.6cm},
]

%% Row 1: byte layout
\node[startcode] (sc) at (0,0) {Start code\\(4 B)};
\node[header, right=0pt of sc] (nh) {NAL header\\(1 B)};
\node[header, right=0pt of nh] (hdr) {Header\\bytes 1–18\\(18 B)};
\node[dc, right=0pt of hdr] (dc) {\textbf{DC bytes}\\offset 19–43\\(25 B)};
\node[rest, right=0pt of dc] (rest) {Remaining\\payload\\($\ldots$)};

%% Annotations below
\draw[<->, thick, orange!70!black]
  ([yshift=-0.15cm]dc.south west) -- ([yshift=-0.15cm]dc.south east)
  node[midway, below, font=\footnotesize] {\texttt{DC\_OFFSET=19, DC\_SIZE=25}};

\draw[dashed, gray]
  (sc.north west) -- ++(0, 0.5) node[above, font=\tiny] {nal\_pos};
\draw[dashed, orange!70!black]
  (dc.north west) -- ++(0, 0.5) node[above, font=\tiny] {nal\_pos+19};
\draw[dashed, orange!70!black]
  (dc.north east) -- ++(0, 0.5) node[above, font=\tiny] {nal\_pos+43};

%% Label
\node[above=0.8cm of hdr, font=\small\itshape] {NAL Unit VCL (Type 1 / Type 5)};

\end{tikzpicture}

\caption{Bố cục byte trong NALU VCL: vùng DC (bytes 19--43) được tô màu cam;
         bytes 0--18 là header không mã hoá}
\label{fig:dc_offset}
\end{figure}

\textbf{Lưu ý quan trọng:} DC\_OFFSET = 19 được xác định bằng phân tích thực nghiệm
trên file cụ thể này (H.264 Baseline Profile). Với các profile khác (Main, High),
vị trí này có thể thay đổi --- đây là hạn chế được thảo luận trong Chương 7.

Các bytes 0--18 (header NALU) chứa thông tin slice như: frame number, picture order
count, reference frame index --- không được mã hoá để đảm bảo khả năng decode
frame (với error concealment) ngay cả khi DC bị mã hoá.

\begin{figure}[ht]
\centering
%% TikZ: H.264 NALU structure — bitsteram level
\begin{tikzpicture}[
  box/.style={draw, minimum height=0.75cm, text centered, font=\small},
  stream/.style={box, fill=gray!10, minimum width=1.2cm},
  nal/.style={box, fill=blue!12},
  vcl/.style={box, fill=green!15},
  nonvcl/.style={box, fill=yellow!20},
  arrow/.style={->, >=stealth},
]

%% Bitstream row
\node[stream, minimum width=0.8cm] (s0) at (0,0) {\ldots};
\node[nal, right=0pt of s0, minimum width=1.4cm] (sc1) {SC\\(4B)};
\node[nal, right=0pt of sc1, minimum width=0.8cm] (t7) {SPS\\T=7};
\node[nal, right=0pt of t7, minimum width=1.4cm] (sc2) {SC\\(4B)};
\node[nal, right=0pt of sc2, minimum width=0.8cm] (t8) {PPS\\T=8};
\node[nal, right=0pt of t8, minimum width=1.4cm] (sc3) {SC\\(4B)};
\node[vcl, right=0pt of sc3, minimum width=1.2cm] (t5) {IDR\\T=5};
\node[nal, right=0pt of t5, minimum width=1.4cm] (sc4) {SC\\(4B)};
\node[vcl, right=0pt of sc4, minimum width=1.2cm] (t1a) {P\\T=1};
\node[nal, right=0pt of t1a, minimum width=1.4cm] (sc5) {SC\\(4B)};
\node[vcl, right=0pt of sc5, minimum width=1.2cm] (t1b) {P\\T=1};
\node[stream, right=0pt of t1b, minimum width=0.8cm] (s1) {\ldots};

%% Zoom in to IDR NALU
\node[vcl, below=1.8cm of t5, minimum width=1.2cm] (idt) {IDR (T=5)};
\node[box, right=0pt of idt, fill=gray!20, minimum width=1.4cm] (hdr18) {Header\\1–18};
\node[box, right=0pt of hdr18, fill=orange!40, minimum width=2.0cm] (dc25) {\textbf{DC}\\25 bytes};
\node[box, right=0pt of dc25, fill=blue!10, minimum width=1.6cm] (rem) {Remain\\payload};

%% Arrow from T5 down
\draw[arrow, dashed] (t5.south) -- (idt.north);

%% DC annotation
\draw[<->, orange!80!black, thick]
  ([yshift=-0.15cm]dc25.south west) -- ([yshift=-0.15cm]dc25.south east)
  node[midway, below, font=\footnotesize] {offset 19..43};

%% Legend
\node[vcl, minimum width=0.6cm, minimum height=0.4cm] (leg1) at (0,-3.4) {};
\node[right=2pt of leg1, font=\footnotesize] {VCL (Type 1, 5) — được mã hoá};
\node[nonvcl, minimum width=0.6cm, minimum height=0.4cm, right=2.8cm of leg1] (leg2) {};
\node[right=2pt of leg2, font=\footnotesize] {Non-VCL (SPS, PPS) — không mã hoá};

\end{tikzpicture}

\caption{Cấu trúc luồng bit H.264: SPS, PPS, IDR (Type 5) và P-frames (Type 1);
         zoom vào IDR NALU cho thấy vùng DC 25 bytes}
\label{fig:nalu_structure}
\end{figure}

\subsection{Mã hoá chọn lọc (Selective Encryption)}

\subsubsection{Nguyên lý và phân loại}

Mã hoá chọn lọc \cite{cheng_selective} là cách tiếp cận mã hoá dựa trên tính
quan trọng tri giác (perceptual importance) của các thành phần video. Thay vì
mã hoá toàn bộ bitstream (chi phí cao), chỉ mã hoá phần tối thiểu đủ để:
\begin{enumerate}
  \item Làm mờ nội dung với người xem không có khóa (không thể hiểu nội dung).
  \item Duy trì khả năng truyền tải (bitstream có thể được parse/route mà không cần decrypt).
  \item Giảm chi phí tính toán xuống phần nhỏ so với full encryption.
\end{enumerate}

Các mức độ mã hoá chọn lọc từ ít đến nhiều:
\begin{itemize}
  \item \textbf{Header-only:} Chỉ mã hoá SPS/PPS --- file không decode được nhưng
        thiếu an toàn nếu attacker đoán được header.
  \item \textbf{DC-only (cách tiếp cận của chúng tôi):} Mã hoá hệ số DC --- balance
        tốt giữa chi phí và hiệu quả bảo vệ.
  \item \textbf{DC+Motion Vector:} Kết hợp DC và motion vectors để bảo vệ tốt hơn
        cho P/B-frames.
  \item \textbf{DC+AC low frequency:} Mã hoá DC và một số AC coefficients tần số thấp
        --- chi phí tăng nhưng bảo vệ mạnh hơn.
\end{itemize}

\subsubsection{Hai chiến lược S1 và S2}

Trong nghiên cứu này, chúng tôi so sánh hai chiến lược mã hoá DC:

\textbf{S1 --- Mã hoá toàn bộ VCL:}
Mã hoá DC của tất cả NALU có Type 1 (P/B-frame) và Type 5 (I-frame).
Với file test: 19.037 NALU được mã hoá / 19.348 tổng = 98.4\% NALU.
S1 bảo vệ mạnh nhất vì cả I-frame lẫn P/B-frames đều bị nhiễu loạn.

\textbf{S2 --- Chỉ mã hoá I-frame:}
Chỉ mã hoá DC của NALU Type 5 (IDR slice / I-frame).
Với file test: 155 NALU được mã hoá / 19.348 tổng = 0.8\% NALU.
S2 hy sinh bảo vệ P/B-frames để đổi lấy chi phí tính toán thấp hơn 120 lần.

\textbf{Lý do I-frame quan trọng hơn:} I-frame là keyframe độc lập --- toàn bộ
frame được mã hoá intra, không tham chiếu frame khác. Khi I-frame bị nhiễu loạn,
lỗi có thể lan truyền sang P/B-frames phụ thuộc vào nó. Tuy nhiên, trong thực tế,
H.264 có cơ chế error concealment nên P/B-frames vẫn hiển thị được (có méo hình).

\subsection{Lý thuyết hỗn độn (Chaos Theory) trong mật mã học}

\subsubsection{Tính chất của hệ hỗn độn}

Các hệ hỗn độn (chaotic systems) có ba tính chất cốt lõi làm cho chúng phù hợp
cho mật mã học \cite{chen_chaos_video}:

\begin{enumerate}
  \item \textbf{Nhạy cảm với điều kiện đầu (Sensitivity to Initial Conditions):}
        Hai quỹ đạo bắt đầu gần nhau sẽ phân kỳ theo hàm mũ theo thời gian.
        Được định lượng bởi \textbf{Lyapunov exponent} $\lambda > 0$: hai điểm
        ban đầu cách nhau $\delta_0$ sẽ cách nhau $\delta_0 e^{\lambda t}$ sau
        thời gian $t$. Với mật mã: thay đổi 1 bit khóa dẫn đến keystream
        hoàn toàn khác.

  \item \textbf{Đặc tính hỗn loạn (Topological Mixing):} Quỹ đạo hỗn độn
        ``trộn'' đồng đều không gian trạng thái --- không có vùng nào bị bỏ sót.
        Điều này đảm bảo tính phân phối đều của keystream.

  \item \textbf{Tính xác định (Deterministic):} Mặc dù có vẻ ngẫu nhiên,
        hệ hỗn độn hoàn toàn xác định: cùng điều kiện đầu → cùng quỹ đạo.
        Điều này cho phép giải mã: sinh lại đúng keystream từ cùng khóa.
\end{enumerate}

Sự tương đồng giữa hệ hỗn độn và mật mã học:

\begin{center}
\begin{tabular}{ll}
\toprule
Mật mã học & Hệ hỗn độn \\
\midrule
Khóa K & Điều kiện đầu $x_0$ \\
Keystream & Quỹ đạo hỗn độn $\{x_n\}$ \\
Độ nhạy khóa & Lyapunov exponent $\lambda$ \\
Pseudorandomness & Topological mixing \\
\bottomrule
\end{tabular}
\end{center}

\subsubsection{PLCM --- Piecewise Linear Chaotic Map}

PLCM (Piecewise Linear Chaotic Map) \cite{li_plcm} là ánh xạ hỗn độn 1D trên
đoạn $[0, 1)$ được định nghĩa bởi:

\begin{equation}
f(x, p) = \begin{cases}
  x / p                    & \text{nếu } 0 \leq x < p \\
  (x - p)/(0.5 - p)        & \text{nếu } p \leq x < 0.5 \\
  f(1 - x, p)              & \text{nếu } 0.5 \leq x < 1
\end{cases}
\label{eq:plcm}
\end{equation}

Với tham số điều khiển $p \in (0, 0.5)$. Đây là ánh xạ tuyến tính từng đoạn
(piecewise linear) với hai nhánh: nhánh tăng dần và nhánh giảm dần (đối xứng).

\textbf{Tính chất toán học của PLCM:}
\begin{itemize}
  \item \textbf{Lyapunov exponent:} $\lambda = -p\log_2 p - (0.5 - p)\log_2(0.5 - p)$
        luôn dương với mọi $p \in (0, 0.5)$, đảm bảo hành vi hỗn độn.
        Giá trị $\lambda$ đạt cực đại tại $p = 0.25$ ($\lambda = 1$).

  \item \textbf{Phân phối đều:} Với $p \in (0, 0.5)$, PLCM có \textbf{invariant measure}
        là phân phối đều trên $[0, 1)$, nghĩa là keystream sinh ra có phân phối
        đồng đều --- tốt cho mật mã.

  \item \textbf{So sánh với Logistic Map:} Logistic Map $x_{n+1} = rx_n(1-x_n)$
        phổ biến hơn nhưng có phân phối phi tuyến. PLCM ưu việt hơn về
        tính phân phối đều và dễ tính toán (phép tính tuyến tính).
\end{itemize}

\textbf{Giới hạn IEEE 754 double precision:}
Trên máy tính thực tế, $x_0$ được biểu diễn bởi số dấu phẩy động double precision
(64 bit, 52 bit mantissa). Do đó keyspace thực tế của PLCM bị giới hạn bởi
$\approx 2^{53}$ giá trị phân biệt trong $(0.0001, 0.4999)$ --- đây là hạn chế
quan trọng của thiết kế hiện tại.

\subsection{Arnold 2D Cat Map}

\subsubsection{Định nghĩa và ma trận biến đổi}

Arnold Cat Map \cite{arnold_catmap} là biến đổi tuyến tính modular trên lưới
$N \times N$ được định nghĩa bởi:

\begin{equation}
\begin{pmatrix} x' \\ y' \end{pmatrix}
= \mathbf{A} \begin{pmatrix} x \\ y \end{pmatrix} \pmod{N}
= \begin{pmatrix} 1 & P \\ Q & PQ+1 \end{pmatrix}
\begin{pmatrix} x \\ y \end{pmatrix} \pmod{N}
\label{eq:arnold}
\end{equation}

với $P, Q$ là các số nguyên dương và $N$ là kích thước lưới. Ma trận $\mathbf{A}$
có $\det(\mathbf{A}) = (PQ+1) - PQ = 1$, nên biến đổi là khả nghịch (invertible)
và bảo toàn diện tích (area-preserving).

\textbf{Tham số trong nghiên cứu này:} $P = 3$, $Q = 11$, $N = 5$ (lưới 5$\times$5,
tương ứng với 25 bytes DC).

Ma trận biến đổi:
\begin{equation}
\mathbf{A} = \begin{pmatrix} 1 & 3 \\ 11 & 34 \end{pmatrix}
\end{equation}

\subsubsection{Áp dụng cho mảng 1D}

Mảng DC 25 bytes được reshape thành ma trận $5 \times 5$:
\begin{itemize}
  \item Byte $i$ tương ứng tọa độ $(x, y) = (i \bmod 5,\; \lfloor i/5 \rfloor)$.
  \item Sau Arnold map: tọa độ mới $(x', y') = \mathbf{A}(x, y) \pmod{5}$.
  \item Chỉ số mới: $i' = y' \times 5 + x'$.
  \item Hoán vị: $\text{permuted}[i'] = \text{DC}[i]$.
\end{itemize}

\subsubsection{Tính chất bảo mật}

\begin{itemize}
  \item \textbf{Tính tuần hoàn (Periodicity):} Arnold map là biến đổi tuần hoàn
        trên $\mathbb{Z}_N^2$. Sau một số lần lặp hữu hạn, lưới trở về trạng thái ban đầu.
        Với $N = 5$, chu kỳ là 20 (mọi điểm trở về vị trí ban đầu sau 20 lần áp dụng Arnold map).
        Trong nghiên cứu này chỉ dùng 1 lần lặp nên tránh được vấn đề tuần hoàn.

  \item \textbf{Số hoán vị phân biệt:} Với $N = 5$ và $P, Q \in \{1,\ldots,4\}$,
        tổng số ma trận Arnold khác nhau $\approx 4 \times 4 = 16 = 2^4$.
        Cộng với tham số iterations (giả sử 1--4): $\approx 2^5$ keyspace.

  \item \textbf{Mục đích chính:} Arnold map không phải nguồn entropy chính,
        mà là lớp hoán vị để tăng diffusion: sau hoán vị, các bytes liền kề
        trong DC không còn liền kề về mặt byte index, làm giảm tương quan cục bộ.
\end{itemize}

\subsubsection{Biến đổi ngược (Inverse Arnold)}

Ma trận nghịch đảo của $\mathbf{A}$ modulo $N$:
\begin{equation}
\mathbf{A}^{-1} = \begin{pmatrix} PQ+1 & -P \\ -Q & 1 \end{pmatrix} \pmod{N}
= \begin{pmatrix} 34 & -3 \\ -11 & 1 \end{pmatrix} \pmod{5}
= \begin{pmatrix} 4 & 2 \\ 4 & 1 \end{pmatrix}
\end{equation}

Hoán vị ngược: $\text{DC}[i] = \text{permuted}[i']$ với $(x, y) = \mathbf{A}^{-1}(x', y') \pmod{5}$.

\subsection{Các thuật toán mã hoá cổ điển}

\subsubsection{AES-128-CTR (Advanced Encryption Standard)}

AES \cite{fips_aes} (FIPS PUB 197) là block cipher được NIST chuẩn hoá năm 2001.
AES-128 dùng khóa 128-bit, xử lý các block 128-bit (16 bytes) qua 10 vòng
(rounds) với các phép biến đổi: SubBytes (S-box substitution), ShiftRows, MixColumns,
và AddRoundKey.

\textbf{Chế độ CTR (Counter Mode):}
CTR biến AES thành stream cipher bằng cách mã hoá bộ đếm tuần tự:
\begin{equation}
ks_i = \text{AES}_K(\text{nonce} \| \text{counter}_i), \quad C_i = P_i \oplus ks_i
\end{equation}

Với 25 bytes DC, cần 2 block AES: block 1 cho bytes 0--15, block 2 cho bytes 16--24.

\textbf{Ưu điểm:} Tốc độ cao nhờ AES-NI hardware instruction trên x86-64;
keystream song song hoá được (parallel keystream generation).

\textbf{Điểm yếu trong ngữ cảnh này:} Khi counter = 0 cố định cho mọi NALU,
keystream $ks = \text{AES}_K(0 \| 0)$ là hằng số --- identitcal keystream reuse
cho mọi NALU (phân tích chi tiết trong Chương 4).

\subsubsection{DES-OFB (Data Encryption Standard --- Output Feedback Mode)}

DES \cite{fips_des} (FIPS PUB 46-3) là block cipher 64-bit với khóa 64-bit (56-bit
hiệu dụng sau khi loại bỏ 8 parity bits), được thiết kế năm 1977. DES có 16 vòng
với Feistel network structure. Hiện nay DES không còn an toàn (brute-force attack
khả thi với phần cứng hiện đại) nhưng được giữ lại trong nghiên cứu này để so sánh.

\textbf{Chế độ OFB (Output Feedback Mode):}
\begin{equation}
O_i = \text{DES}_K(O_{i-1}), \quad O_0 = IV, \quad C_i = P_i \oplus O_i
\end{equation}

OFB tạo keystream độc lập với plaintext, cho phép mã hoá trước (pre-computation).
Khi IV = 0 cố định, mọi NALU nhận cùng keystream --- điểm yếu tương tự AES-CTR.

\textbf{So sánh với AES-CTR:} DES chậm hơn AES $\approx$ 8 lần trên x86-64 vì
thiếu hardware acceleration và S-box lookup table (64-bit→4-bit) không tối ưu.

\subsubsection{RC4 (Rivest Cipher 4)}

RC4 \cite{rc4} là stream cipher được thiết kế bởi Ron Rivest năm 1987, ban đầu
là trade secret của RSA Security, sau bị leak năm 1994.

\textbf{Cấu trúc RC4:}
\begin{enumerate}
  \item \textbf{KSA (Key Scheduling Algorithm):} Khởi tạo S-box 256 byte:
        $S = [0, 1, \ldots, 255]$, sau đó hoán vị dựa trên khóa K.
        \begin{equation}
        j = (j + S[i] + K[i \bmod \text{keylen}]) \bmod 256; \quad S[i] \leftrightarrow S[j]
        \end{equation}

  \item \textbf{PRGA (Pseudo-Random Generation Algorithm):} Sinh keystream byte-by-byte:
        \begin{equation}
        i = (i+1)\bmod 256;\; j = (j+S[i])\bmod 256;\; S[i]\leftrightarrow S[j];\; k = S[(S[i]+S[j])\bmod 256]
        \end{equation}
\end{enumerate}

\textbf{Điểm yếu thống kê:} Fluhrer-Mantin-Shamir attack \cite{fms_rc4} chứng minh
rằng PRGA output ban đầu (đặc biệt là byte đầu tiên) có thiên lệch thống kê
(statistical bias) --- xác suất xuất hiện 0 là $2/256$ thay vì $1/256$ lý thuyết.

Trong nghiên cứu này, RC4 được khởi tạo lại từ đầu cho mỗi NALU với cùng khóa K,
nên mọi NALU nhận cùng 25 bytes PRGA output --- điểm yếu cốt lõi.

\subsection{Các chỉ số đánh giá bảo mật}

\subsubsection{NPCR và UACI --- Độ nhạy plaintext}

Wu và cộng sự \cite{wu_npcr} đề xuất NPCR (Number of Pixels Change Rate) và
UACI (Unified Average Changing Intensity) để đo \textbf{độ nhạy plaintext}:
khả năng thay đổi 1 bit plaintext gây ra thay đổi ciphertext trên diện rộng.

Cho $P$ là DC bytes gốc và $P'$ là phiên bản với byte đầu flip 1 bit:
$P'[0] = P[0] \oplus 1$, $P'[i] = P[i]$ với $i > 0$.

Gọi $C_1 = \text{Enc}(P, K)$ và $C_2 = \text{Enc}(P', K)$. Định nghĩa:

\begin{equation}
D[i] = \begin{cases} 0 & \text{nếu } C_1[i] = C_2[i] \\ 1 & \text{nếu } C_1[i] \neq C_2[i] \end{cases}
\end{equation}

\begin{equation}
\text{NPCR} = \frac{\sum_{i=0}^{24} D[i]}{25} \times 100\%
\label{eq:npcr}
\end{equation}

\begin{equation}
\text{UACI} = \frac{1}{25} \sum_{i=0}^{24} \frac{|C_1[i] - C_2[i]|}{255} \times 100\%
\label{eq:uaci}
\end{equation}

\textbf{Ngưỡng lý thuyết:}
\begin{itemize}
  \item NPCR: Nếu mỗi byte ciphertext thay đổi ngẫu nhiên độc lập khi 1 bit
        plaintext thay đổi, $\text{NPCR}_{\text{exp}} = (1 - 2^{-8}) \times 100\% \approx 99.61\%$.
        Ngưỡng thực tế: NPCR $> 99.6\%$.

  \item UACI: Nếu $|C_1[i] - C_2[i]|$ phân phối đều trên $\{0, \ldots, 255\}$,
        $\text{UACI}_{\text{exp}} = \frac{1}{256} \sum_{d=0}^{255} \frac{d}{255} \times 100\%
        = \frac{255/2}{255} \times 100\% = 50\%$. Tuy nhiên, với điều kiện
        byte-difference $\neq 0$ (khi thay đổi xảy ra): $\text{UACI}_{\text{exp}} \approx 33.46\%$.
\end{itemize}

\subsubsection{Entropy Shannon}

Entropy Shannon \cite{shannon1949} đo độ ngẫu nhiên (randomness) của chuỗi byte:

\begin{equation}
H = -\sum_{v=0}^{255} p(v) \log_2 p(v) \quad \text{(bits/byte)}
\label{eq:entropy}
\end{equation}

trong đó $p(v)$ là xác suất xuất hiện của giá trị byte $v$ trong dữ liệu.
Entropy tối đa là $H = 8$ bits/byte (phân phối đều trên 256 giá trị).
Ngưỡng thực tế: $H > 7.9$ bits/byte.

Trong nghiên cứu này, Entropy được tính trên tất cả DC bytes encrypted từ tất cả
NALU (không phải từng NALU riêng lẻ). Kết quả $H$ cao chứng tỏ DC bytes sau mã hoá
có phân phối gần đều --- dấu hiệu keystream tốt.

\subsubsection{SSIM --- Structural Similarity Index}

SSIM \cite{wang_ssim} đo tương đồng cấu trúc giữa hai tín hiệu, xét đồng thời
luminance (độ sáng), contrast (độ tương phản), và structure (cấu trúc):

\begin{equation}
\text{SSIM}(x, y) = \frac{(2\mu_x\mu_y + c_1)(2\sigma_{xy} + c_2)}{(\mu_x^2 + \mu_y^2 + c_1)(\sigma_x^2 + \sigma_y^2 + c_2)}
\end{equation}

SSIM $\in [-1, 1]$, với 1 là giống hệt, 0 là không tương quan, $-1$ là phản tương quan.

Trong nghiên cứu này, SSIM được tính trực tiếp trên DC bytes (coi như tín hiệu 1D
25 phần tử) giữa phiên bản gốc và phiên bản encrypted. SSIM DC $\approx 0$ là tốt.

\subsubsection{TSSIM --- Temporal SSIM}

TSSIM (Temporal SSIM) đo tương đồng cấu trúc theo chiều thời gian giữa các frames
liên tiếp. Với DC bytes của frame $t$ và frame $t+1$:
\begin{equation}
\text{TSSIM}(t) = \text{SSIM}\bigl(\text{DC}(t), \text{DC}(t+1)\bigr)
\end{equation}

TSSIM baseline (video gốc không mã hoá) phản ánh sự liên tục tự nhiên giữa các frames.
Sau mã hoá, TSSIM tốt là $<$ baseline: mã hoá phá vỡ cấu trúc temporal.

\subsubsection{PSNR --- Peak Signal-to-Noise Ratio}

PSNR đo chất lượng tái tạo tín hiệu theo thang logarithm:

\begin{equation}
\text{PSNR} = 10 \log_{10} \frac{\text{MAX}^2}{\text{MSE}}, \quad
\text{MSE} = \frac{1}{N}\sum_{i=0}^{N-1}(x_i - y_i)^2
\end{equation}

Với DC bytes (byte values $\in [0, 255]$), MAX = 255.
\begin{itemize}
  \item \textbf{PSNR DC encrypted:} Đo giữa DC gốc và DC encrypted. Giá trị thấp
        ($< 10$ dB) là tốt --- chứng tỏ mã hoá gây ra nhiễu loạn lớn.
  \item \textbf{PSNR original$\to$decrypted:} Đo giữa file gốc và file giải mã.
        $\text{PSNR} = \infty$ (MSE = 0) là bắt buộc --- giải mã hoàn toàn chính xác.
\end{itemize}

\subsubsection{BER và SHA-256 --- Tính đúng đắn}

\textbf{BER (Bit Error Rate):}
\begin{equation}
\text{BER} = \frac{\text{số bit khác nhau giữa gốc và giải mã}}{\text{tổng số bit}}
\end{equation}

BER = 0 là yêu cầu bắt buộc. Bất kỳ bit lỗi nào cũng có thể lan truyền và gây
decode error cho cả GOP (Group of Pictures).

\textbf{SHA-256:} Hash cryptographic toàn bộ file ($> 219$ MB) giữa gốc và giải mã.
SHA-256 PASS là kiểm tra byte-exact mạnh nhất: xác suất collision $\approx 2^{-256}$.

\subsection{Tổng kết chỉ số đánh giá}

\begin{table}[ht]
\centering
\caption{Tổng kết các chỉ số đánh giá và ngưỡng mục tiêu}
\label{tab:metrics_summary}
\begin{tabular}{llll}
\toprule
Nhóm & Chỉ số & Ngưỡng tốt & Ghi chú \\
\midrule
\multirow{3}{*}{Bảo mật} & NPCR (\%) & $> 99.6$ & Chỉ hybrid \\
 & UACI (\%) & $\approx 33.46$ & Chỉ hybrid \\
 & Entropy (bits/byte) & $> 7.9$ & Tất cả algo \\
\midrule
\multirow{3}{*}{Chất lượng} & PSNR DC (dB) & $< 10$ & Thấp = tốt \\
 & SSIM DC & $\approx 0$ & Gần 0 = tốt \\
 & TSSIM DC & $<$ baseline & Thấp hơn gốc \\
\midrule
\multirow{2}{*}{Đúng đắn} & BER & $= 0$ & Bắt buộc \\
 & SHA-256 & PASS & Bắt buộc \\
\midrule
Hiệu năng & DC time (ns/NALU) & Càng thấp càng tốt & So sánh tương đối \\
\bottomrule
\end{tabular}
\end{table}
